% Define colours for code listings
\definecolor{keyword}{RGB}{0,0,128}
\definecolor{comment}{RGB}{0,96,0}
\definecolor{string}{RGB}{128,0,0}
\lstset{
  language=Python,
  basicstyle=\ttfamily\small,
  keywordstyle=\color{keyword}\bfseries,
  commentstyle=\color{comment}\itshape,
  stringstyle=\color{string},
  showstringspaces=false,
  frame=single,
  framerule=0.5pt,
  numbers=left,
  numberstyle=\tiny,
  numbersep=5pt
}


\chapter{Lab 3: Scaling the Customer‑Support Agent with AgentCore Gateway and Identity}
\label{chap:Lab 3: Scaling the Customer‑Support Agent with AgentCore Gateway and Identity}

\section{Introduction and Motivation}

The previous labs equipped you with a locally running customer‑support agent
complete with custom tools and persistent conversation memory.  However, the
prototype still suffers from limitations common to early experiments: each
agent bundles its own copy of every tool, there is no centralised security
control and your laptop cannot sustainably serve multiple users.  Lab 3
addresses these shortcomings by introducing the \emph{AgentCore Gateway} and
\emph{AgentCore Identity}.  Together they allow tools to be published once
and reused across any number of agents, while ensuring that only authorised
callers can invoke them.

The Gateway is a managed server that exposes functions (e.g., AWS Lambda
functions, REST APIs or other services) as \emph{Model Context Protocol
\,MCP\,} tools.  MCP standardises how tools are described and invoked so
that any agent framework can discover and use them.  AgentCore Identity
provides a unified authentication layer: inbound authentication verifies
that the caller has a valid OAuth token (issued by identity providers
such as Cognito, Okta or Entra), while outbound authentication securely
retrieves credentials needed to call downstream services.  These
capabilities transform your agent into an enterprise‑ready solution
without rewriting your existing tool logic.

\section{Prerequisites}

Before beginning this lab, ensure the following requirements are met:

\begin{itemize}
  \item \textbf{Lab 2 completed.}  You should already have a working
    customer‑support agent with memory integration.  The memory resource
    identifier (\texttt{memory\_id}) and session management code will be reused.
  \item \textbf{AWS credentials configured.}  The Gateway and Identity
    services require permissions to create resources such as Cognito user
    pools, IAM roles and Lambda functions.  It is assumed that the
    accompanying CloudFormation template has provisioned a user pool, a
    gateway IAM role, a Lambda execution role and DynamoDB tables for
    warranty and customer profile data.
  \item \textbf{Python environment.}  Continue using the virtual
    environment from previous labs and ensure that the following packages
    are installed:

\begin{lstlisting}[language=bash]
pip install strands-agents bedrock-agentcore-starter-toolkit mcp-client boto3
\end{lstlisting}

  \item \textbf{External tools.}  Two Lambda functions have been provided:
    \texttt{check\_warranty\_status} to query warranty information from a
    DynamoDB table and \texttt{web\_search} to perform DuckDuckGo searches.
    These functions are defined in the \texttt{prerequisite/lambda/python}
    folder of the workshop repository.  We will expose them through the
    Gateway.
\end{itemize}

\section{Importing Required Libraries}

As with previous labs, begin by importing the necessary Python modules and
initialising clients.  The AgentCore Gateway is controlled via the
\texttt{bedrock-agentcore-control} client; the \texttt{MCPClient} from the
\texttt{strands} library helps us call Gateway tools from within the
agent.

\begin{lstlisting}[language=Python]
# Import standard libraries
import os
import sys
import boto3

# Strands Agent and model classes
from strands import Agent
from strands.models import BedrockModel

# MCP client for connecting to Gateway tools
from strands.tools.mcp import MCPClient
from mcp.client.streamable_http import streamablehttp_client

# Helper functions provided by the workshop repository
from lab_helpers.utils import (
    get_or_create_cognito_pool,
    put_ssm_parameter,
    get_ssm_parameter,
    load_api_spec,
)

# Instantiate the regional clients
REGION = boto3.session.Session().region_name
gateway_client = boto3.client(
    "bedrock-agentcore-control",
    region_name=REGION,
)

print("✅ Libraries imported successfully!")

\end{lstlisting}

This snippet obtains the current AWS region and creates a client for the
AgentCore control plane.  It also imports a utility function
\texttt{get\_or\_create\_cognito\_pool} that will create or retrieve a
Cognito user pool for authentication.  The \texttt{load\_api\_spec}
function loads the API specification file that describes the tools.

\section{Exposing Existing Functions as MCP Tools}

The first step in centralising tools is to migrate functions from local
modules into the Gateway.  In previous labs the warranty and web search
capabilities were implemented directly within the agent.  Here they will be
run inside AWS Lambda functions.  A thin wrapper in the Lambda handler
inspects the incoming event to determine which tool to call.  The event
contains the tool name (passed via the custom context field
\texttt{bedrockAgentCoreToolName}) and a JSON payload with user inputs.

An excerpt of the Lambda handler illustrates the routing logic:

\begin{lstlisting}[language=Python]
def lambda_handler(event, context):
    """Dispatches calls to the appropriate tool based on the tool name.

    Parameters
    ----------
    event: dict
        The event contains 'body' with JSON-serialised tool parameters and
        a custom context entry 'bedrockAgentCoreToolName' specifying which
        tool to invoke.
    context: object
        Lambda context providing runtime information.

    Returns
    -------
    dict
        A response dictionary with a status code and JSON body.
    """
    # Extract the tool name from the Lambda context
    extended_tool_name = context.client_context.custom["bedrockAgentCoreToolName"]
    tool_name = extended_tool_name.split("___")[1]

    # Dispatch to the appropriate function
    if tool_name == "check_warranty_status":
        serial_number = event["body"].get("serial_number")
        customer_email = event["body"].get("customer_email")
        result = check_warranty_status(serial_number, customer_email)
        return {"statusCode": 200, "body": result}
    elif tool_name == "web_search":
        keywords = event["body"].get("keywords")
        region = event["body"].get("region", "us-en")
        max_results = event["body"].get("max_results", 5)
        result = web_search(keywords, region, max_results)
        return {"statusCode": 200, "body": result}
    else:
        return {"statusCode": 400, "body": f"Unknown tool: {tool_name}"}

\end{lstlisting}

The helper functions \texttt{check\_warranty\_status} and \texttt{web\_search}
contain the actual logic.  The warranty function looks up a serial number
and (optionally) verifies ownership via the customer email; the web search
function uses the DuckDuckGo library to return a list of search results.

\section{Defining Tool Schemas with an API Specification}

The Gateway requires a specification file that describes each tool’s name,
description and input schema.  This file uses a JSON array where each
object corresponds to one tool.  Below is the specification for the
\texttt{check\_warranty\_status} and \texttt{web\_search} tools.  Note
that \texttt{serial\_number} is mandatory for the warranty check, while
\texttt{keywords} is mandatory for the web search.  Optional fields have
defaults.

\begin{lstlisting}[language=json]
[
  {
    "name": "check_warranty_status",
    "description": "Check the warranty status of a product using its serial number and optionally verify via email",
    "inputSchema": {
      "type": "object",
      "properties": {
        "serial_number": { "type": "string" },
        "customer_email": { "type": "string" }
      },
      "required": ["serial_number"]
    }
  },
  {
    "name": "web_search",
    "description": "Search the web for updated information using DuckDuckGo",
    "inputSchema": {
      "type": "object",
      "properties": {
        "keywords": {
          "type": "string",
          "description": "The search query keywords"
        },
        "region": {
          "type": "string",
          "description": "The search region (e.g., us-en, uk-en, ru-ru)"
        },
        "max_results": {
          "type": "integer",
          "description": "The maximum number of results to return"
        }
      },
      "required": ["keywords"]
    }
  }
]
\end{lstlisting}

This specification lives at \texttt{prerequisite/lambda/api\_spec.json} in
the workshop repository.  It will be passed to the Gateway when creating
the Lambda target in Step \ref{sec:add-target}.

\section{Creating the AgentCore Gateway with OAuth Authentication}

To expose the Lambda functions via MCP, you must create an AgentCore
Gateway.  The Gateway unifies all tools under a single URL and enforces
security via inbound authentication.  In this example we use a custom
JWT authorizer that relies on Amazon Cognito to validate OAuth access
tokens.  The discovery URL and client ID are retrieved from Systems
Manager Parameter Store by the helper function
\texttt{get\_or\_create\_cognito\_pool}.  The Gateway itself is created
using the \texttt{create\_gateway} call on the control client.

\begin{lstlisting}[language=Python]
# Name of the Gateway
gateway_name = "customersupport-gw"

# Retrieve or create a Cognito user pool and client
cognito_config = get_or_create_cognito_pool(refresh_token=True)

# Define custom JWT authorizer configuration
auth_config = {
    "customJWTAuthorizer": {
        "allowedClients": [cognito_config["client_id"]],
        "discoveryUrl": cognito_config["discovery_url"],
    }
}

try:
    print(f"Creating gateway in region {REGION} with name: {gateway_name}")
    response = gateway_client.create_gateway(
        name=gateway_name,
        roleArn=get_ssm_parameter("/app/customersupport/agentcore/gateway_iam_role"),
        protocolType="MCP",
        authorizerType="CUSTOM_JWT",
        authorizerConfiguration=auth_config,
        description="Customer Support AgentCore Gateway",
    )

    # Persist gateway identifiers for later use
    gateway_id = response["gatewayId"]
    put_ssm_parameter("/app/customersupport/agentcore/gateway_id", gateway_id)
    put_ssm_parameter("/app/customersupport/agentcore/gateway_name", gateway_name)
    put_ssm_parameter("/app/customersupport/agentcore/gateway_arn", response["gatewayArn"])
    put_ssm_parameter("/app/customersupport/agentcore/gateway_url", response["gatewayUrl"])
    print(f"✅ Gateway created successfully with ID: {gateway_id}")
except Exception:
    # If gateway already exists, retrieve its details from Parameter Store
    gateway_id = get_ssm_parameter("/app/customersupport/agentcore/gateway_id")
    response = gateway_client.get_gateway(gatewayIdentifier=gateway_id)
    print(f"Found existing gateway with ID: {gateway_id}")

# Consolidate gateway details in a dictionary for later use
gateway = {
    "id": gateway_id,
    "name": gateway_name,
    "gateway_url": get_ssm_parameter("/app/customersupport/agentcore/gateway_url"),
    "gateway_arn": get_ssm_parameter("/app/customersupport/agentcore/gateway_arn"),
}
\end{lstlisting}

The call to \texttt{create\_gateway} registers a new MCP server.  It
specifies the ARN of an IAM role that allows the Gateway to invoke
the Lambda functions and access Parameter Store.  If a gateway already
exists, the code simply fetches its identifier and continues.

\section{Attaching the Lambda Target to the Gateway\label{sec:add-target}}

With the Gateway created, the next step is to attach the Lambda functions
as a \emph{target}.  A target bundles the backend implementation (our
Lambda function), the tool specification and credentials used to call the
backend.  Gateway supports multiple targets, enabling you to group related
tools together or connect to different microservices under a single
endpoint.

\begin{lstlisting}[language=Python]
try:
    api_spec_file = "./prerequisite/lambda/api_spec.json"
    if not os.path.exists(api_spec_file):
        print(f"❌ API specification file not found: {api_spec_file}")
        sys.exit(1)

    # Load the tool definitions
    api_spec = load_api_spec(api_spec_file)

    # Configure the target to invoke our Lambda function via MCP
    lambda_target_config = {
        "mcp": {
            "lambda": {
                "lambdaArn": get_ssm_parameter("/app/customersupport/agentcore/lambda_arn"),
                "toolSchema": {"inlinePayload": api_spec},
            }
        }
    }

    # Use the Gateway IAM role to supply credentials to the Lambda backend
    credential_config = [{"credentialProviderType": "GATEWAY_IAM_ROLE"}]

    # Create the target on the Gateway
    response = gateway_client.create_gateway_target(
        gatewayIdentifier=gateway_id,
        name="LambdaUsingSDK",
        description="Lambda Target using SDK",
        targetConfiguration=lambda_target_config,
        credentialProviderConfigurations=credential_config,
    )

    print(f"✅ Gateway target created: {response['targetId']}")
except Exception as e:
    print(f"❌ Error creating gateway target: {str(e)}")
\end{lstlisting}

This code loads the API specification, constructs a target configuration
specifying the Lambda ARN and tool schema, and sets the credential provider
type to \texttt{GATEWAY\_IAM\_ROLE}.  The IAM role defined earlier grants
the Gateway permission to execute the Lambda.  On success, the target
identifier is printed.

\section{Building a Secure MCP Client and Agent}

Once the Gateway and target are configured, you can create a secure client
to call the tools.  The \texttt{MCPClient} from Strands accepts a factory
function that returns a streaming HTTP client connected to the Gateway URL
with a valid Authorization header.  The Cognito helper returns a bearer
token for the currently authenticated user.

\subsection{Set up the MCP Client}

\begin{lstlisting}[language=Python]
# Display the Gateway URL for reference
print(f"Gateway Endpoint - MCP URL: {gateway['gateway_url']}")

# Create the MCP client using a bearer token from Cognito
mcp_client = MCPClient(
    lambda: streamablehttp_client(
        gateway["gateway_url"],
        headers={"Authorization": f"Bearer {cognito_config['bearer_token']}"},
    )
)
\end{lstlisting}

The \texttt{streamablehttp\_client} is provided by the MCP library.  It
establishes an HTTP connection to the Gateway and automatically streams
responses.  Passing a lambda into \texttt{MCPClient} defers client
creation until the context manager is entered, ensuring that tokens are
fresh.

\subsection{Create the Agent with Local and Gateway Tools}

The final step is to assemble the agent, combining both local tools and
those retrieved from the Gateway.  The following function constructs a
Strands \texttt{Agent} using the Haiku model from Bedrock, the memory
configuration from Lab 2 and the list of tools.  It must be called inside
a \texttt{with mcp\_client} context so that tool discovery happens with a
valid connection.

\begin{lstlisting}[language=Python]
from lab_helpers.lab1_strands_agent import (
    get_product_info,
    get_return_policy,
    get_technical_support,
    SYSTEM_PROMPT,
)

import uuid
from lab_helpers.lab2_memory import create_or_get_memory_resource
from bedrock_agentcore.memory.integrations.strands.config import AgentCoreMemoryConfig, RetrievalConfig
from bedrock_agentcore.memory.integrations.strands.session_manager import AgentCoreMemorySessionManager

# Create or retrieve the memory resource
memory_id = create_or_get_memory_resource()

# Generate a new session and define the customer identifier
SESSION_ID = str(uuid.uuid4())
CUSTOMER_ID = "customer_001"

# Configure memory retrieval for preferences and semantic memories
memory_config = AgentCoreMemoryConfig(
    memory_id=memory_id,
    session_id=SESSION_ID,
    actor_id=CUSTOMER_ID,
    retrieval_config={
        "support/customer/{actorId}/semantic": RetrievalConfig(top_k=3, relevance_score=0.2),
        "support/customer/{actorId}/preferences": RetrievalConfig(top_k=3, relevance_score=0.2),
    },
)

# Initialise the Bedrock model
model = BedrockModel(
    model_id="global.anthropic.claude-haiku-4-5-20251001-v1:0",
    temperature=0.3,
    region_name=REGION,
)

def create_agent(prompt: str) -> str:
    """Constructs and invokes a customer support agent using local and MCP tools.

    Parameters
    ----------
    prompt : str
        The user's natural language question.

    Returns
    -------
    str
        The agent's textual response.
    """
    try:
        with mcp_client:
            # Combine local tools with those exposed via the Gateway
            tools = [
                get_product_info,
                get_return_policy,
                get_technical_support,
            ] + mcp_client.list_tools_sync()

            # Create the agent
            agent = Agent(
                model=model,
                tools=tools,
                system_prompt=SYSTEM_PROMPT,
                session_manager=AgentCoreMemorySessionManager(memory_config, REGION),
            )

            # Invoke the agent and return the response
            return agent(prompt)
    except Exception as e:
        raise RuntimeError(f"Agent invocation failed: {str(e)}")

print("✅ Customer support agent created successfully!")
\end{lstlisting}

The function calls \texttt{mcp\_client.list\_tools\_sync()} to obtain a list
of tool descriptors from the Gateway.  These descriptors are callable
objects; when the agent selects one, the MCP client automatically
marshals the request into a JSON payload, adds the OAuth token and sends
it to the Gateway.  The agent continues to use local tools for
product information, return policy and knowledge base retrieval, since
these remain specific to the customer support use case.

\section{Testing the Enhanced Agent}

To verify that the gateway integration works properly, run a series of
test prompts.  The following code iterates through several examples and
prints the agent’s response.  You should see both local and MCP tools
being used: for instance, queries about warranty status will trigger the
\texttt{check\_warranty\_status} tool via the Gateway, while troubleshooting
questions will call the knowledge‑base tool.

\begin{lstlisting}[language=Python]
test_prompts = [
    "List all of your tools",  # Should list local and Gateway tools
    "I bought an iPhone 14 last month. It heats up—how do I fix it?",
    "I have a Gaming Console Pro device. I want to check my warranty status. The serial number is MNO33333333.",
    "What are the warranty support guidelines?",
    "How can I fix a Lenovo Thinkpad with a blue screen?",
    "Tell me detailed information about the technical documentation on installing a new CPU",
]

def test_agent_responses(prompts):
    for i, prompt in enumerate(prompts, 1):
        print(f"\nTest Case {i}: {prompt}")
        print("-" * 50)
        try:
            response = create_agent(prompt)
            print(response)
        except Exception as e:
            print(f"Error: {str(e)}")
        print("-" * 50)

# Run the tests
test_agent_responses(test_prompts)

print("\n✅ Basic testing completed!")
\end{lstlisting}

Running this function should confirm that the agent now recognises five
tools: the three local tools and the two MCP tools from the Gateway.  A
query about warranty will produce a structured response from the Lambda
function, while general troubleshooting questions will still leverage the
knowledge base.  You may notice slightly longer latency when the agent
calls remote tools due to network overhead.

\section{Summary and Next Steps}

In this lab you transformed a simple, locally bound agent into a modular
system capable of securely accessing centralised tools.  The AgentCore
Gateway provides a single entry point for multiple tools, enabling
reusability across different agents and use cases.  AgentCore Identity
ensures that only authorised callers can invoke these tools and that
downstream services are accessed with the correct credentials.  By
delegating common capabilities (like warranty checks and web searches) to
shared infrastructure, you eliminate duplication and simplify maintenance.

You have also seen how to integrate the Gateway with an existing agent:
initialising an MCP client with a bearer token, combining local and
remote tools, and configuring memory to maintain state across sessions.
The testing script demonstrates the agent’s ability to seamlessly invoke
both local and centralised tools.

Lab 3 leaves you with a functional yet still locally executed prototype.
In Lab 4 you will deploy this agent to the serverless AgentCore Runtime,
gain comprehensive observability and enable automatic scaling.  You will
also address remaining limitations such as manual resource management and
lack of production‑ready infrastructure.

